\section{Тест производительности}
{\itshape В данном разделе проводится сравнение производительности реализованного красно-чёрного дерева с контейнером стандартной библиотеки C++ \texttt{std::map}, который также основан на красно-чёрном дереве.}

Тест производительности представляет собой сравнение времени выполнения основных операций словаря: вставки, поиска и удаления. Для каждой структуры данных выполняется одинаковый набор операций над одинаковыми входными данными.

В качестве входных данных используются случайно сгенерированные строки длины 10 символов, состоящие из букв английского алфавита. Генерация производится с фиксированным зерном генератора случайных чисел, что обеспечивает одинаковые наборы данных для всех тестируемых структур.

Для каждого теста выполняются следующие шаги:
\begin{itemize}
    \item генерация $N$ случайных ключей;
    \item вставка всех ключей в структуру;
    \item поиск всех ключей;
    \item удаление всех ключей.
\end{itemize}

Время измеряется с использованием библиотеки \texttt{<chrono>} и включает только выполнение операций над структурами данных (без учёта ввода-вывода).

Тестирование проводилось для трёх размеров входных данных:
\begin{itemize}
    \item $10^4$ элементов;
    \item $10^5$ элементов;
    \item $10^6$ элементов.
\end{itemize}

\begin{alltt}
root@container:/workspaces/lab# g++ banch.cpp -Wall -o banch
root@container:/workspaces/lab# ./banch 10000

=== Comparison for 10000 operations ===

RBTree:
  Insert: 13.822 ms
  Find:   10.759 ms
  Remove: 10.646 ms

std::map:
  Insert: 8.741 ms
  Find:   6.368 ms
  Remove: 8.781 ms

root@container:/workspaces/lab# ./banch 100000

=== Comparison for 100000 operations ===

RBTree:
  Insert: 162.333 ms
  Find:   154.704 ms
  Remove: 154.954 ms

std::map:
  Insert: 105.560 ms
  Find:   84.695 ms
  Remove: 116.283 ms

root@container:/workspaces/lab# ./banch 1000000

=== Comparison for 1000000 operations ===

RBTree:
  Insert: 2630.753 ms
  Find:   2530.334 ms
  Remove: 2582.752 ms

std::map:
  Insert: 1536.105 ms
  Find:   1349.625 ms
  Remove: 1867.914 ms
\end{alltt}

Как видно из результатов, контейнер \texttt{std::map} стабильно показывает более высокую производительность по сравнению с реализованным красно-чёрным деревом.

Это объясняется несколькими причинами:
\begin{itemize}
    \item стандартная библиотека использует высоко оптимизированную реализацию;
    \item возможны дополнительные оптимизации компилятора и более эффективная работа с памятью;
    \item в моей реализации присутствуют дополнительные накладные расходы (например, рекурсивные вызовы и менее оптимальная балансировка).
\end{itemize}

При этом обе реализации демонстрируют асимптотическую сложность $O(\log n)$ для всех операций, что подтверждается ростом времени выполнения при увеличении размера входных данных.

Таким образом, несмотря на более низкую практическую производительность, реализованное красно-чёрное дерево корректно работает и соответствует теоретическим ожиданиям по сложности операций.

\pagebreak